\documentclass[a4paper,10pt]{article}

\usepackage[utf8]{inputenc}
\usepackage[T1]{fontenc}
\usepackage[french]{babel}
\usepackage{amsmath,amssymb}
\usepackage{geometry}
\usepackage{tikz}
\usepackage{tcolorbox}
\usepackage{multicol}
\usepackage{enumitem}
\geometry{margin=1.2cm}
\pagestyle{empty}

\begin{document}

% ======================= PAGE 1 =======================

\begin{center}
    {\LARGE \textbf{Python}}
\end{center}

\begin{multicols}{3}

\begin{tcolorbox}[title=Conditions]
if condition :

\hspace*{5mm}...

elif ... :

\hspace*{5mm}...

else :

\hspace*{5mm}...
\end{tcolorbox}

\begin{tcolorbox}[title=Fonctions]
\texttt{def f(a,b=0,c=0):}

\hspace*{5mm}instruction

\hspace*{5mm}return ...

\vspace{2mm}

\texttt{type(var)} retourne le type de la variable.
\end{tcolorbox}

\begin{tcolorbox}[title=Boucles]
while ... :

\hspace*{5mm}...

\vspace{2mm}

for x in sequence :

\hspace*{5mm}...

\vspace{2mm}

for x in range(2,7)

range(7) $\rightarrow$ 0 à 6

range(2,7,2)

$\rightarrow$ 2 à 6 step by 2

\vspace{2mm}

enumerate(list)

$\rightarrow$ sequence of tuples

(index, element)
\end{tcolorbox}

\columnbreak

\begin{tcolorbox}[title=Division]
/ division

// division entière

\% reste
\end{tcolorbox}

\begin{tcolorbox}[title=Booléens]
and

or
\end{tcolorbox}

\begin{tcolorbox}[title=Chaînes de caractères]
"simple quote"

'double quote'

'''triple double quote'''

"""docs"""
\end{tcolorbox}

\begin{tcolorbox}[title=Lists]
L = [ ]

L = [1, ``txt'']

\vspace{2mm}

list.c = [1]*10

list(range(5))

0 à 4

\vspace{2mm}
\begin{itemize}[leftmargin=5mm]
    \item len(list)
    \item min(list)
    \item max(list)
\end{itemize}

if ``...'' in list:

\hspace*{5mm}...

else :

\hspace*{5mm}...

\vspace{2mm}

\textbf{Méthodes :}

\begin{itemize}[leftmargin=5mm]
    \item len(list)
    \item list.append(x)
    \item list.insert(i,x)
    \item list.extend(list2)
    \item list.pop()
    \item list.pop(i)
    \item list.remove(x)
\end{itemize}
\end{tcolorbox}

\columnbreak

\begin{tcolorbox}[title=Extraction d'éléments]
list[i:start:j-end excluded]

list[i:j:end excluded]

list[i:start:]

list[i:j:step-size]

\vspace{2mm}

On peut aussi faire :

\texttt{a = ["1",2,"("]}
\end{tcolorbox}

\begin{tcolorbox}[title=Dictionnaires]
dict = \{ ... : ... \}

\vspace{2mm}

dict["..."] = ...

\vspace{2mm}
.keys()\quad .values()

.pop(key)
\end{tcolorbox}

\begin{tcolorbox}[title=Tuples]
Similaire à lists mais fixed in size and immutable.

\vspace{2mm}

tuple = (...,...)

\vspace{2mm}

On ne peut pas modifier.

Je peux typer une list sans typecast en tuple.
\end{tcolorbox}

\begin{tcolorbox}[title=Fonctions utiles]
obj = max(...)

return ..., ...

\vspace{2mm}

a,b = min,max(...)
\end{tcolorbox}

\end{multicols}

\newpage

% ======================= PAGE 2 =======================

\begin{multicols}{2}

\begin{tcolorbox}[title=List comprehension]
[thing for thing in list]

\vspace{2mm}

Ex :

listA = [1,2,3,4]

listB = [2*elem for elem in listA]

\vspace{2mm}

[thing for thing in list if condition]
\end{tcolorbox}

\begin{tcolorbox}[title=NumPy]
import numpy as np

\vspace{2mm}

\textbf{Creating arrays}

Convert normal python arrays :

np.array([...])

\vspace{2mm}

Using built-in functions :

np.zeros((2,3))

np.ones(...)

np.eye(4)

np.arange(5,11)

np. ones like(...)

\vspace{2mm}

np.array([1,0], dtype=np.int32)

astype(np.float32)

\vspace{2mm}

.ndim

.size

.shape

.dtype

\vspace{2mm}

Get elements :

x[2::2, ::2]

\vspace{2mm}

np.concatenate([arr1,arr2], axis=0)

np.stack()

np.sum()

np.min \quad np.max \quad np.argmin

np.random.rand(3,2)
\end{tcolorbox}

\columnbreak

\begin{tcolorbox}[title=Matplotlib]
import numpy

import numpy as np

import matplotlib.pyplot as plt

\vspace{2mm}

fig = plt.figure() $\rightarrow$ creates a figure

\vspace{2mm}

plt.plot(x,y, marker='*', color='r')

plt.xlabel("x")

plt.ylabel("y")

plt.title("title")

plt.show() $\rightarrow$ pour afficher

plt.close(fig) $\rightarrow$ close the plot

\vspace{2mm}

On peut avoir plusieurs fonctions sur le graph

en faisant plusieurs plt.plot(...)

\vspace{2mm}

plt.scatter(...) $\rightarrow$ même chose que plot

mais sans les lignes.

\vspace{2mm}

img = plt.imread("img/...jpg")

$\rightarrow$ reads an image

plt.imshow(img)

\vspace{2mm}

Masking :

mask = x > 20

x[mask]
\end{tcolorbox}

\end{multicols}

\end{document}
